\chapter{Markov chains and the AR(1) laboratory}
\label{ch:markov-chains}

The paper's data are sequences drawn from a first-order Markov chain with a linear
autoregressive transition and a non-Gaussian innovation. The choice is not
decorative: it is the one family in which the covariance is fixed by two numbers and
every non-Gaussian feature lives in the cumulants of order four and higher, so that
``second-order structure'' and ``higher-order structure'' are separated by
construction and not by an approximation. This chapter defines Markov chains and
transition kernels, distinguishes strict from second-order stationarity (the paper's
chain has the second and, for non-Gaussian innovations, not the first), derives the
covariance $\Cov(a_i,a_j)=\corr^{|i-j|}$ for every innovation law, studies the
covariance matrix $\Sig_0$ and its spectrum, derives the chain's fourth-cumulant tensor,
and fixes what the two sizes $\nsites$ and $\nseq$ mean.

\section{The Markov property and transition kernels}

\begin{definition}[Markov chain]
A sequence of real random variables $a_1,a_2,\dots,a_{\nsites}$ is a \term{Markov
chain} if, for every $i$, the conditional law of $a_i$ given $(a_1,\dots,a_{i-1})$
depends only on $a_{i-1}$:
\begin{equation}
  \Prob(a_i\in B\mid a_1,\dots,a_{i-1})=\Prob(a_i\in B\mid a_{i-1})\qquad\text{for every Borel set }B.
\end{equation}
\end{definition}

When the conditional laws have densities, the chain is described by an
\term{initial density} $\mu$ and \term{transition kernels} $\kernel_i(a'\mid a)$, the
density of $a_i$ at $a'$ given $a_{i-1}=a$. The kernel is a density in its first
argument for every value of the second, $\int\kernel_i(a'\mid a)\dd a'=1$. Multiplying
the conditional densities along the sequence gives the joint density
\begin{equation}
  p(a_1,\dots,a_{\nsites})=\mu(a_1)\prod_{i=2}^{\nsites}\kernel_i(a_i\mid a_{i-1}),
  \label{eq:markov-joint}
\end{equation}
and conversely any density of this product form is Markov, because conditioning on
$(a_1,\dots,a_{i-1})$ leaves a function of $a_i$ that depends on the past only through
the factor $\kernel_i(a_i\mid a_{i-1})$. The chain is \term{time-homogeneous} if every
$\kernel_i$ is the same kernel $\kernel$; the paper's chain is. The pair $(\mu,\kernel)$
then determines the law of the whole sequence, which is why ``learning the
dynamics'' in the paper means learning $\kernel$ (the initial law $\mu=\normal01$ is
held fixed and known).

\begin{figure}[t]
\centering
\begin{tikzpicture}[>=Latex,node distance=1.5cm]
  \tikzset{site/.style={circle,draw,minimum size=0.8cm,inner sep=0pt}}
  \node[site] (a1) {$a_1$};
  \node[site,right=of a1] (a2) {$a_2$};
  \node[site,right=of a2] (a3) {$a_3$};
  \node[right=of a3] (dots) {$\cdots$};
  \node[site,right=of dots] (aL) {$a_{\nsites}$};
  \draw[->] (a1) -- node[above] {$\kernel$} (a2);
  \draw[->] (a2) -- node[above] {$\kernel$} (a3);
  \draw[->] (a3) -- node[above] {$\kernel$} (dots);
  \draw[->] (dots) -- node[above] {$\kernel$} (aL);
  \node[above=0.5cm of a1] {$\mu$};
  \node[below=0.35cm of a2] {$\varepsilon_2$};
  \node[below=0.35cm of a3] {$\varepsilon_3$};
  \node[below=0.35cm of aL] {$\varepsilon_{\nsites}$};
\end{tikzpicture}
\caption{The chain: the initial law $\mu$ at site $1$, one transition kernel $\kernel$
per edge, and one independent innovation $\varepsilon_i$ entering at each site $i\ge2$.
The posterior factor graph of the paper's Figure 1 is this chain with one observation
attached to each site.}
\label{fig:chain}
\end{figure}

Two consequences of \cref{eq:markov-joint} are used constantly. The
\term{Chapman--Kolmogorov equation} says that the two-step kernel is the composition
$\kernel^{(2)}(a''\mid a)=\int\kernel(a''\mid a')\kernel(a'\mid a)\dd a'$, and more
generally the $k$-step kernel is the $k$-fold composition. And conditioning on a site
splits the chain: given $a_i$, the past $(a_1,\dots,a_{i-1})$ and the future
$(a_{i+1},\dots,a_{\nsites})$ are independent, which is what makes the posterior factor
graph of the paper a tree and the forward--backward recursion exact (chapter on
factor graphs and belief propagation).

\section{Stationarity: strict and second-order}

\begin{definition}[Invariant law and stationarity]
A density $\pi$ is \term{invariant} for $\kernel$ if $\int\kernel(a'\mid a)\pi(a)\dd a=\pi(a')$
for every $a'$. A time-homogeneous chain started from an invariant law,
$\mu=\pi$, is \term{strictly stationary}: the joint law of $(a_i,\dots,a_{i+k})$ does
not depend on $i$. A sequence is \term{second-order stationary} if its means are
constant and $\Cov(a_i,a_j)$ depends only on $|i-j|$.
\end{definition}

Strict stationarity implies second-order stationarity when second moments exist; the
converse fails, and the paper's chain is the example. Its first site is $\normal01$,
which is the invariant law of the transition only when the innovation is Gaussian;
for every other family the chain is second-order stationary (next section) but not
strictly stationary, and its fourth cumulant relaxes toward the invariant value along
the chain, as the chapter on moments and cumulants computed. The paper's proof of
Proposition 1 says explicitly that ``nothing assumes stationarity'', and its
fourth-cumulant tensor (\cref{sec:kappa4-chain}) carries the site dependence exactly.
The word ``stationary'' in this compendium always says which sense is meant.

\section{The AR(1) chain}

\begin{definition}[First-order autoregression]
Let $\innov$ be a density on $\Reals$ with mean zero and variance $\ivar$, and let
$|\corr|<1$. The \term{AR(1) chain} with innovation $\innov$ is
\begin{equation}
  a_1\sim\mu,\qquad a_i=\corr a_{i-1}+\varepsilon_i\quad(i\ge2),\qquad
  \varepsilon_2,\varepsilon_3,\dots\ \text{i.i.d.\ with density }\innov,\ \text{independent of }a_1 .
  \label{eq:ar1}
\end{equation}
Its transition kernel is $\kernel(a'\mid a)=\innov(a'-\corr a)$.
\end{definition}

The kernel is a function of the residual $a'-\corr a$ alone, which is what makes the
grid representation of the chapter on grid quadrature a matrix of residuals; the
helper \code{\_outer\_residual} in \code{src/markov\_score/priors.py} builds that
matrix once for every family. Because
$\varepsilon_i$ is independent of $(a_1,\dots,a_{i-1})$, the law of $a_i$ given the
whole past depends only on $a_{i-1}$: the recursion is Markov.

\paragraph{The invariant law.}
Iterating \cref{eq:ar1}, $a_i=\corr^{i-1}a_1+\sum_{m=2}^i\corr^{i-m}\varepsilon_m$. As
$i\to\infty$ the contribution of $a_1$ vanishes and the sum converges to the
moving-average series $\sum_{k\ge0}\corr^k\varepsilon_{-k}$ of the chapter on moments
and cumulants; its law is the invariant law of the kernel, whatever $\mu$ was. For
Gaussian innovations it is $\normal{0}{\ivar/(1-\corr^2)}$, which at $\ivar=1-\corr^2$ is
$\normal01$: the paper's Gaussian chain is strictly stationary. For a Laplace
innovation the invariant law is the law of $\sum_k\corr^k\varepsilon_{-k}$, which has no
closed form but has all the cumulants of the chapter on moments and cumulants; it is
not the law of $a_1$.

\section{Unit variance and geometric covariance for every innovation law}

\begin{proposition}[Second-order structure of the chain]
\label{prop:second-order}
Let $\mu=\normal01$ and $\ivar=1-\corr^2$ in \cref{eq:ar1}. Then, for every innovation
density of mean zero and variance $\ivar$,
\begin{equation}
  \E[a_i]=0,\qquad \Var(a_i)=1,\qquad \Cov(a_i,a_j)=\corr^{|i-j|}\quad\text{for all }i,j .
  \label{eq:covariance}
\end{equation}
\end{proposition}
\begin{proof}
Means: $\E a_1=0$, and inductively $\E a_i=\corr\E a_{i-1}+\E\varepsilon_i=0$.
Variances: $\Var(a_1)=1$, and inductively, since $\varepsilon_i$ is independent of
$a_{i-1}$,
\[
  \Var(a_i)=\corr^2\Var(a_{i-1})+\Var(\varepsilon_i)=\corr^2+(1-\corr^2)=1 .
\]
Covariances: for $i>j$, write $a_i=\corr^{i-j}a_j+\sum_{m=j+1}^{i}\corr^{i-m}\varepsilon_m$,
where the sum is independent of $a_j$ and has mean zero, so
$\Cov(a_i,a_j)=\corr^{i-j}\Var(a_j)=\corr^{i-j}$. Only the mean and variance of $\innov$
were used.
\end{proof}

\why{This is the sentence ``the variance choice makes $\Var(a_i)=1$ and
$\Cov(a_i,a_j)=\corr^{|i-j|}$ exactly for every innovation law'' of the paper's Section
2, and it is the whole reason the laboratory separates orders. Every innovation family
of the paper is symmetric, so its odd cumulants vanish; \cref{prop:second-order} pins
the cumulants of order two; therefore two chains with different innovation families
share every cumulant tensor of order one, two and three, and differ only from order
four on. A Gaussian-innovation kernel fitted by maximum likelihood can at best recover
$(\corr,\ivar)$, hence $\Sig_0$; whatever it then leaves on the table is, by this
proposition, fourth order or higher. That is what the paper's Proposition 1 turns
into a computable floor.}

The code parameterises every family by $\corr$ alone, computing $\ivar=1-\corr^2$ and
the family's scale from it (\cref{tab:sampling}); the theory module
\code{src/markov\_score/theory.py} (function \code{chain\_covariance}) keeps $\ivar$
free, because the two-parameter family $(\corr,\ivar)$ is what the Fisher information
of the chapter on Fisher information is computed on. For $\ivar\ne1-\corr^2$ the
induction gives $\Var(a_i)=\corr^{2(i-1)}+\ivar(1-\corr^{2(i-1)})/(1-\corr^2)$, drifting from
$1$ toward $\ivar/(1-\corr^2)$, and $\Cov(a_i,a_j)=\corr^{|i-j|}\Var(a_{\min(i,j)})$; this
is what makes $\partial_\ivar\Sig_0\ne0$, so that $\ivar$ is a second-order coordinate
in the paper's sense.

\section{The covariance matrix and its spectrum}
\label{sec:spectrum}

At $\ivar=1-\corr^2$ the covariance matrix of the chain is the symmetric Toeplitz matrix
\begin{equation}
  \Sig_0=\bigl[\corr^{|i-j|}\bigr]_{i,j=1}^{\nsites}
  =\begin{pmatrix}1&\corr&\corr^2&\cdots&\corr^{\nsites-1}\\ \corr&1&\corr&\cdots&\corr^{\nsites-2}\\
  \vdots&&\ddots&&\vdots\\ \corr^{\nsites-1}&\corr^{\nsites-2}&\cdots&\corr&1\end{pmatrix},
  \label{eq:kms}
\end{equation}
known as the \term{Kac--Murdock--Szeg\H{o} matrix}. Three of its properties matter.

\paragraph{Determinant and inverse.}
By the chain rule for densities, $\det\Sig_0=\Var(a_1)\prod_{i\ge2}\Var(a_i\mid a_{i-1})=1\cdot\ivar^{\nsites-1}=(1-\corr^2)^{\nsites-1}$.
The inverse (the \term{precision matrix}) is tridiagonal,
\begin{equation}
  \Sig_0^{-1}=\frac1{\ivar}\begin{pmatrix}1&-\corr&&&\\ -\corr&1+\corr^2&-\corr&&\\ &\ddots&\ddots&\ddots&\\
  &&-\corr&1+\corr^2&-\corr\\ &&&-\corr&1\end{pmatrix},
  \label{eq:precision}
\end{equation}
which one verifies by multiplying \cref{eq:kms} by \cref{eq:precision} row by row
(the interior row $i$ of the product has entries $\corr^{|i-j|}(-\corr\cdot\corr^{\mp1}+(1+\corr^2)-\corr\cdot\corr^{\pm1})/\ivar$
for $j\ne i$, which vanish, and $(1-\corr^2)/\ivar=1$ on the diagonal). A banded
precision matrix is the linear-algebraic face of the Markov property: for a Gaussian
vector, $(\Sig^{-1})_{ij}=0$ means $a_i$ and $a_j$ are conditionally independent given
the rest. This is the matrix \code{GaussianLongRange.precision} in
\code{src/markov\_score/priors.py} builds before perturbing it.

\paragraph{Eigenvalue bounds.}
\begin{proposition}[Spectrum of the KMS matrix]
\label{prop:kms-spectrum}
Every eigenvalue $\lambda$ of $\Sig_0$ satisfies
\begin{equation}
  \frac{1-|\corr|}{1+|\corr|}\le\lambda\le\frac{1+|\corr|}{1-|\corr|},
  \label{eq:kms-bounds}
\end{equation}
and $\lmax(\Sig_0)$ increases with $\nsites$ to the upper bound.
\end{proposition}
\begin{proof}
Upper bound. For a symmetric matrix the largest eigenvalue is at most the largest
absolute row sum (the operator norm is bounded by the $\infty$-norm), and every row sum
of \cref{eq:kms} is at most $1+2\sum_{k\ge1}|\corr|^k=(1+|\corr|)/(1-|\corr|)$. Lower
bound. $\Sig_0$ is the covariance of the vector $\va$, so for any unit vector $\vect v$,
$\vect v^{\!\top}\Sig_0\vect v=\Var(\vect v^{\!\top}\va)$. Writing $\va$ in terms of the
innovations, $\vect v^{\!\top}\va=\sum_mc_m\xi_m$ with $\xi_1=a_1$, $\xi_m=\varepsilon_m$
independent, and one checks from the moving-average coefficients that this variance
is the quadratic form of the Toeplitz matrix with symbol $g(\omega)=(1-\corr^2)/(1-2\corr\cos\omega+\corr^2)$;
a symmetric Toeplitz form is bounded below by the minimum of its symbol, which is
$(1-\corr^2)/(1+|\corr|)^2=(1-|\corr|)/(1+|\corr|)$. The monotonicity in $\nsites$ holds
because the $\nsites\times\nsites$ matrix is a principal submatrix of the
$(\nsites+1)\times(\nsites+1)$ one (Cauchy interlacing), and the limit is the maximum of
the symbol, $g(0)=(1+\corr)/(1-\corr)$ for $\corr>0$, by the Szeg\H{o} theory of Toeplitz
forms.
\end{proof}

\paragraph{The expansion parameter and the crossover.}
The paper's asymptotic theorems expand in $e^{-2t}$, but for the Gaussian chain the
natural parameter is $e^{-2t}\lmax(\Sig_0)$: the noised covariance
$\Sig_t=e^{-2t}\Sig_0+\Dt\Id=\Id+e^{-2t}(\Sig_0-\Id)$ is a perturbation of the identity
whose size in operator norm is $e^{-2t}(\lmax-1)$. The \term{crossover time}
\begin{equation}
  \tstar=\tfrac12\log\lmax(\Sig_0)\le\tfrac12\log\frac{1+\corr}{1-\corr}=\operatorname{artanh}\corr
  \label{eq:tstar}
\end{equation}
is the time at which $e^{-2t}\lmax=1$; below it the longest-wavelength covariance mode
is still of order one and the asymptotic rates are not yet visible, above it they
are. The bound $\operatorname{artanh}\corr$ is exact as $\nsites\to\infty$; at $\nsites=32$
the paper's value, a generated macro, is a little smaller. The paper uses
$\tstar$ as an exact scale for the Gaussian chain (where the chapter on Fisher
information shows how it enters the closed-form information) and as a heuristic for
the non-Gaussian ones, and reads every ``shallower than predicted'' slope of the
earlier literature as a measurement below $\tstar$. The neighbouring time
$\tfrac12\log(\lmax-1)$, at which $\Pt/\gamma$ enters $L^2(\gamma)$ (chapter on
Gaussian integrals and Hermite polynomials), is smaller by a little; the two coincide
in spirit and the paper names only $\tstar$.

\section{The fourth-cumulant tensor of the chain}
\label{sec:kappa4-chain}

\begin{proposition}[Fourth-cumulant tensor]
\label{prop:kappa4-chain}
For the chain \cref{eq:ar1} with $\mu=\normal01$ and an innovation with fourth cumulant
$\kappa_4(\varepsilon)$, for sites $i,j,k,l\ge1$,
\begin{equation}
  \kappa_4[ijkl]=\kappa_4(\varepsilon)\sum_{m=2}^{\min(i,j,k,l)}\corr^{\,i+j+k+l-4m},
  \label{eq:kappa4-chain}
\end{equation}
with the empty sum equal to zero (any index equal to $1$).
\end{proposition}
\begin{proof}
Write each site as a linear combination of the independent variables
$\xi_1=a_1,\ \xi_m=\varepsilon_m$ ($m\ge2$): $a_i=\sum_{m=1}^ic_{im}\xi_m$ with
$c_{i1}=\corr^{i-1}$ and $c_{im}=\corr^{i-m}$ for $2\le m\le i$. By the transformation rule
and additivity of cumulant tensors (chapter on moments and cumulants), and because the
$\xi_m$ are independent so that only ``diagonal'' joint cumulants
$\operatorname{cum}(\xi_m,\xi_m,\xi_m,\xi_m)=\kappa_4(\xi_m)$ survive,
\[
  \kappa_4[ijkl]=\sum_mc_{im}c_{jm}c_{km}c_{lm}\,\kappa_4(\xi_m).
\]
The Gaussian $\xi_1$ has $\kappa_4=0$; each innovation contributes
$\kappa_4(\varepsilon)\corr^{(i-m)+(j-m)+(k-m)+(l-m)}$ for $m\le\min(i,j,k,l)$, which is
\cref{eq:kappa4-chain}.
\end{proof}

On the diagonal, $\kappa_4[iiii]=\kappa_4(\varepsilon)(1-\corr^{4(i-1)})/(1-\corr^4)$, the
site-dependent value already met, which tends to the stationary
$\kappa_4(\varepsilon)/(1-\corr^4)$. The Frobenius norm $\Frob{\kfour}^2$ grows linearly
in $\nsites$ (each site far from the start contributes the same amount), so the paper's
score-gap constant $\tfrac16\Frob{\kfour}^2$ is proportional to the chain length and
the paper's theorems are stated at fixed $\nsites$. The function \code{chain\_kappa4}
of \code{src/markov\_score/theory.py} evaluates \cref{eq:kappa4-chain} as a geometric
sum, and \code{tests/test\_theory\_hermite.py} checks six of its entries against
brute-force fourth cumulants of sampled Laplace chains, including $\kappa_4[1111]=0$.

\section{Sampling a chain}

Sampling is the recursion itself: draw $a_1\sim\normal01$, then for $i=2,\dots,\nsites$
draw $\varepsilon_i$ from the innovation law with variance $\ivar=1-\corr^2$ and set
$a_i=\corr a_{i-1}+\varepsilon_i$. \Cref{tab:sampling} lists how each family's parameter
is derived from $\ivar$. The samplers are the classes \code{GaussianAR1},
\code{LaplaceAR1}, \code{StudentTAR1}, \code{UniformAR1} and \code{GaussianMixtureAR1}
of \code{src/markov\_score/priors.py}, and the generalised Gaussian lives in
\code{src/markov\_score/shape\_kernel.py}. Each class also exposes a method
\code{log\_transition\_matrix} that evaluates the kernel on the residual matrix of a
grid, indexed child first and parent second.

\begin{table}[t]
\centering
\caption{Variance-$\ivar$ parameterisation of the innovation families, as sampled.}
\label{tab:sampling}
\small
\begin{tabular}{@{}lll@{}}
\toprule
family & draw $\varepsilon$ as & parameter from $\ivar$ \\
\midrule
Gaussian & $\sqrt{\ivar}\,Z$ & --- \\
Laplace & $\mathrm{Laplace}(0,b)$ & $b=\sqrt{\ivar/2}$ \\
Student-$t_\nu$ & $s\,T_\nu$ & $s=\sqrt{\ivar(\nu-2)/\nu}$ \\
uniform & $\mathrm{Uniform}(-h,h)$ & $h=\sqrt{3\ivar}$ \\
mixture, separation $\kappa$ & $\pm m+sZ$, sign fair & $m=\sqrt{\kappa\ivar}$, $s=\sqrt{(1-\kappa)\ivar}$ \\
generalised Gaussian, $\beta$ & density $\propto\exp(-|e/s_\beta|^\beta)$ & $s_\beta=\sqrt{\ivar\,\Gamma(1/\beta)/\Gamma(3/\beta)}$ \\
\bottomrule
\end{tabular}
\end{table}

Two details of the implementation are part of the experimental design. First, random
streams are keyed by strings and integers through \code{rng\_for} in
\code{src/markov\_score/utils.py}, so that the same clean chains and the same noise
draws are used by every method compared at a setting; this is what makes the paper's
comparisons paired (chapter on the statistics of the experiments). Second, the
sampler draws the innovations and the recursion exactly, with no truncation; the
grid on which inference runs truncates the domain to $[-\halfwidth,\halfwidth]$, and a
sampled site beyond it is the mechanism the paper's Appendix B discusses for the
heavy-tailed family.

\section{The chain length $\nsites$ and the number of sequences $\nseq$}

The \term{chain length} $\nsites$ is the number of sites of one sequence and, since a
sequence is a point of $\Reals^{\nsites}$, the dimension of the space in which the
diffusion runs. The frozen value is $\nsites=32$; the length study evaluates kernels
fitted at $32$ on chains of length up to $256$, and the regime sweep uses $32$ and
$128$. The \term{number of sequences} $\nseq$ is the number of independent draws of
the chain in a data set; the sizes are $\nseq\in\{32,128,512,2048\}$ in the
comparisons, and the theory's sample-size statements (Fisher information per
sequence, the local asymptotic minimax bound) are in the limit $\nseq\to\infty$ at
fixed $\nsites$.

A data set therefore has $\nsites\nseq$ scalar observations, but they are not
$\nsites\nseq$ independent numbers: within a sequence the sites are dependent, with
the correlation length of $\corr^{|i-j|}$, and the independent unit is the sequence.
This is why every information statement in the paper is ``per sequence'' and why the
constant $\tfrac16\Frob{\kfour}^2$ carries a factor of order $\nsites$: a longer chain
carries more information about its kernel, in proportion to its length, and the
same estimator applies to any length because the kernel is a function on $\Reals^2$
and inference at length $\nsites'$ is the same recursion. The chapter on conventions
explains why these two numbers are not called $n$ and $d$.

\section{Worked example: the frozen chain at $\corr=0.85$}

Take $\corr=0.85$, so $\ivar=1-0.7225=0.2775$. The first four rows and columns of
$\Sig_0$ are, from \cref{eq:kms},
\[
  \begin{pmatrix}1&0.85&0.7225&0.6141\\0.85&1&0.85&0.7225\\0.7225&0.85&1&0.85\\0.6141&0.7225&0.85&1\end{pmatrix},
\]
with determinant $\ivar^3=0.2775^3\approx0.02137$ for the $4\times4$ block. The precision
entries are $1/\ivar\approx3.604$, $(1+\corr^2)/\ivar\approx6.207$ and $-\corr/\ivar\approx-3.063$;
the interior diagonal $6.207$ is the conditional precision of a site given both
neighbours, and $3.604$ that of an end site given its one neighbour.

The row-sum bound of \cref{prop:kms-spectrum} for the infinite chain is
$(1+\corr)/(1-\corr)=1.85/0.15\approx12.33$, and the lower bound is
$0.15/1.85\approx0.081$; the $4\times4$ block's largest row sum is
$1+0.85+0.7225+0.6141\approx3.19$, so at length $4$ the largest eigenvalue is at most
$3.19$, and it grows toward $12.33$ with the length. The crossover bound is
$\operatorname{artanh}(0.85)=\tfrac12\log12.33\approx1.256$; the paper's $\tstar$ at
$\nsites=32$ lies below this and is reported in its Section 3. Already the $2\times2$
block has eigenvalues $1\pm\corr$, i.e. $1.85$ and $0.15$, so $\lmax>2$ from length $3$
on, and the Gaussian chain's $\Pz/\gamma$ is outside $L^2(\gamma)$ for every length the
paper uses.

For Laplace innovations, $\kappa_4(\varepsilon)=3\ivar^2\approx0.2310$, and
\cref{eq:kappa4-chain} gives on the diagonal $\kappa_4[2222]\approx0.231$,
$\kappa_4[3333]\approx0.231(1+0.522)\approx0.352$, $\kappa_4[4444]\approx0.231(1+0.522+0.2725)\approx0.415$,
and off the diagonal, for instance, $\kappa_4[2233]=\kappa_4(\varepsilon)\corr^{2}\approx0.167$
(only $m=2$ contributes, with exponent $2+2+3+3-8=2$).

\where{{\sloppy The paper's Section 2 (the chain, variance matching, the covariance twin),
Section 3 (the crossover $\tstar$, the bound $\lmax\le(1+\corr)/(1-\corr)$, the
$L^2(\gamma)$ remark), Section 5 (the regime sweep in $\corr$ and $\nsites$, the length
study), Appendix A (the proof of Proposition 1, ``nothing assumes stationarity'') and
Appendix B (the kernel on the grid). Experiments E01, E07, E11, E13, E14, E15 and the
length rows of E03. Code: \code{src/markov\_score/priors.py} (the chain classes and
their samplers), \code{src/markov\_score/shape\_kernel.py} (the generalised-Gaussian
kernel), \code{src/markov\_score/theory.py} (\code{chain\_covariance},
\code{chain\_kappa4}, \code{crossover\_time}), \code{experiments/frozen\_config.py}
(the frozen $\corr$, $\nsites$ and family), \code{tests/test\_theory\_hermite.py}.\par}}

\paragraph{References.}
S. P. Meyn and R. L. Tweedie (2009), \emph{Markov Chains and Stochastic Stability},
2nd edition, Cambridge University Press (Chapters 3 and 10: kernels, invariant
measures).
P. J. Brockwell and R. A. Davis (1991), \emph{Time Series: Theory and Methods},
2nd edition, Springer (Chapter 3: AR processes, stationarity and the moving-average
representation).
M. Kac, W. L. Murdock and G. Szeg\H{o} (1953), ``On the eigen-values of certain Hermitian
forms'', Journal of Rational Mechanics and Analysis 2, 767--800.
U. Grenander and G. Szeg\H{o} (1958), \emph{Toeplitz Forms and Their Applications},
University of California Press (Chapter 5: the spectrum of Toeplitz forms and the
symbol).
W. F. Trench (1999), ``Asymptotic distribution of the spectra of a class of generalized
Kac--Murdock--Szeg\H{o} matrices'', Linear Algebra and its Applications 294, 181--192.
